% !TEX root = ../../../main.tex
\section{Oriented area and volume}
\label{sec:linear-algebra-i:oriented-area-and-volume}

% Sources: October 1 PDF, pp. 3--18; IMG_9207--IMG_9208.
\subsection{Parallelograms and their area}

Let \(V\) be a two-dimensional real vector space. The parallelogram spanned
by \(v_1,v_2\in V\) is
\[
  P(v_1,v_2)=\{av_1+bv_2\mid 0\leq a,b\leq1\}.
\]
\autoref{fig:linear-algebra-i:determinant-parallelogram} shows its two
generating vectors and a point given by such a linear combination.

\begin{marginfigure}
  \centering
  % Supplied October 1 PDF, p. 3, and IMG_9207; schematic reconstruction.
\begin{tikzpicture}[
  x=1cm,y=1cm,font=\footnotesize,line width=0.5pt,
  line cap=round,line join=round,
  vector/.style={-{Latex[length=1.6mm,width=1.1mm]}}
]
  \coordinate (O) at (0,0);
  \coordinate (V1) at (0.8,1.8);
  \coordinate (V2) at (2.6,0);
  \coordinate (S) at ($(V1)+(V2)$);
  \coordinate (P) at ($0.45*(V1)+0.45*(V2)$);
  \fill[black!8] (O) -- (V2) -- (S) -- (V1) -- cycle;
  \draw (V1) -- (S) -- (V2);
  \draw[vector] (O) -- (V1) node[pos=0.65,left=5pt] {\(v_1\)};
  \draw[vector] (O) -- (V2) node[midway,below=5pt] {\(v_2\)};
  \draw[vector,black!60,shorten >=1.6pt] (O) -- (P);
  \fill (P) circle[radius=1.4pt];
  \node[above=5pt] at (P) {\(av_1+bv_2\)};
  \node[below left=4pt] at (O) {\(0\)};
\end{tikzpicture}

  \caption[Parallelogram spanned by two vectors.]{A parallelogram spanned by \(v_1,v_2\), shown schematically.
    The interior point has the form \(av_1+bv_2\).}
  \label{fig:linear-algebra-i:determinant-parallelogram}
\end{marginfigure}

Fix an ordered basis \(E=(e_1,e_2)\) and normalize the area of
\(P(e_1,e_2)\) to be \(1\). If
\[
  v_1=ae_1+he_2,\qquad v_2=we_1,\qquad w>0,
\]
then the width is \(w\), the height is \(|h|\), and the ordinary area is
\[
  \operatorname{area}P(v_1,v_2)=|h|w.
\]
For \(h>0\), this is \(hw\); for \(h<0\), the parallelogram lies below
its base. Both configurations appear in
\autoref{fig:linear-algebra-i:determinant-signed-height}.

\begin{figure}[htbp]
  \centering
  % Supplied October 1 PDF, pp. 4--6; schematic basis-coordinate diagrams.
\begin{tikzpicture}[
  x=1cm,y=1cm,font=\footnotesize,line width=0.5pt,
  line cap=round,line join=round,
  vector/.style={-{Latex[length=1.6mm,width=1.1mm]}},
  guide/.style={black!55,dashed}
]
  \begin{scope}
    \coordinate (O) at (0,0);
    \coordinate (V1) at (0.65,1.5);
    \coordinate (V2) at (2.65,0);
    \coordinate (S) at ($(V1)+(V2)$);
    \coordinate (F) at (S |- O);
    \fill[black!8] (O) -- (V2) -- (S) -- (V1) -- cycle;
    \draw[guide] (V2) -- (F);
    \draw[guide] (F) -- (S) node[midway,right=5pt,text=black] {\(h>0\)};
    \draw (V1) -- (S) -- (V2);
    \draw[vector] (O) -- (V2) node[midway,below=5pt] {\(v_2=we_1\)};
    \draw[vector] (O) -- (V1) node[pos=0.75,left=5pt] {\(v_1\)};
    \draw[vector,black!55] (O) -- (0,0.7) node[pos=0.8,left=5pt] {\(e_2\)};
    \node at (1.65,-0.8) {\(A(v_1,v_2)<0\)};
  \end{scope}
  \begin{scope}[xshift=5.5cm,yshift=1.5cm]
    \coordinate (O) at (0,0);
    \coordinate (V1) at (0.65,-1.5);
    \coordinate (V2) at (2.65,0);
    \coordinate (S) at ($(V1)+(V2)$);
    \coordinate (F) at (S |- O);
    \fill[black!8] (O) -- (V2) -- (S) -- (V1) -- cycle;
    \draw[guide] (V2) -- (F);
    \draw[guide] (F) -- (S) node[midway,right=5pt,text=black] {\(h<0\)};
    \draw (V1) -- (S) -- (V2);
    \draw[vector] (O) -- (V2) node[midway,above=5pt] {\(v_2=we_1\)};
    \draw[vector] (O) -- (V1) node[midway,left=5pt] {\(v_1\)};
    \draw[vector,black!55] (O) -- (0,0.7) node[pos=0.8,left=5pt] {\(e_2\)};
    \node at (1.65,-2.3) {\(A(v_1,v_2)>0\)};
  \end{scope}
\end{tikzpicture}

  \caption[Signed height and oriented area.]{Positive and negative height coefficients relative to the
    ordered basis, shown schematically. Ordinary area is \(|h|w\);
    the oriented area \(A(v_1,v_2)=-hw\) changes sign.}
  \label{fig:linear-algebra-i:determinant-signed-height}
\end{figure}

\subsection{Oriented area and bilinearity}

Choose the orientation for which \(e_2\) lies counterclockwise of \(e_1\),
and normalize \(A(e_1,e_2)=1\). Define the oriented area by
\[
  A(v,w)=
  \begin{cases}
    \operatorname{area}P(v,w) & \text{if \((v,w)\) is positively oriented},\\
    -\operatorname{area}P(v,w) & \text{if \((v,w)\) is negatively oriented},\\
    0 & \text{if \(v,w\) are linearly dependent}.
  \end{cases}
\]
Thus \(A(v,w)\) is positive when \(w\) lies counterclockwise of \(v\),
and negative when \(w\) lies clockwise of \(v\).

Equal inputs are linearly dependent. The parallelogram \(P(u,u)\)
collapses to a line segment and has zero area, so \(A(u,u)=0\) for
every \(u\in V\). In particular,
\[
  A(v,v)=A(w,w)=0.
\]
% The author requested this explanation before the coordinate calculation.

In the signed-height configuration,
\[
  A(v_1,v_2)=-hw.
\]
The ordinary area is \(|A(v_1,v_2)|\).
Under rescaling, ordinary area satisfies
\[
  \begin{aligned}
    \operatorname{area}P(cv,w)&=|c|\operatorname{area}P(v,w),\\
    \operatorname{area}P(v,cw)&=|c|\operatorname{area}P(v,w).
  \end{aligned}
\]
% IMG_9207 first records the absolute-value scaling of ordinary area.
% Its later c-scaling identities and PDF pp. 9--12 concern oriented area.

\begin{definition}[Bilinear function]
  \label{thm:linear-algebra-i:bilinear-function}
  A function \(B:V\times V\to\R\) is \emph{bilinear} if it is linear
  in each input separately, with the other input fixed. Thus, for all
  \(v,v_1,v_2,w,w_1,w_2\in V\) and \(c\in\R\),
  \[
    \begin{aligned}
      B(v_1+v_2,w)&=B(v_1,w)+B(v_2,w),\\
      B(v,w_1+w_2)&=B(v,w_1)+B(v,w_2),\\
      B(cv,w)&=B(v,cw)=cB(v,w).
    \end{aligned}
  \]
\end{definition}
% IMG_9207 gives these three identities, with both scalar identities
% written together. B distinguishes the general definition from area A.

The oriented area function \(A\) is bilinear and satisfies these identities.

\autoref{fig:linear-algebra-i:determinant-area-addition} illustrates
addition with a common base. When the two vectors lie on opposite sides
of the base, the same identity can be written as
\[
  A(v_1,w)=A(v_1+v_2,w)-A(v_2,w).
\]
% PDF p. 8 labels a downward v_2 as having A(v_2,w)<0 while drawing w
% to the right. This conflicts with the orientation convention on pp. 6
% and 11. The valid addition identity is retained without that inequality.

\begin{figure}[htbp]
  \centering
  % Supplied October 1 PDF, pp. 7--8; schematic common-base decomposition.
\begin{tikzpicture}[
  x=1cm,y=1cm,font=\footnotesize,line width=0.5pt,
  line cap=round,line join=round,
  vector/.style={-{Latex[length=1.6mm,width=1.1mm]}}
]
  \begin{scope}
    \coordinate (O) at (0,0);
    \coordinate (W) at (2.4,0);
    \coordinate (V1) at (-0.65,1.05);
    \coordinate (V2) at (0.8,1.05);
    \coordinate (S) at ($(V1)+(V2)$);
    \coordinate (V1W) at ($(V1)+(W)$);
    \coordinate (SW) at ($(S)+(W)$);
    \fill[black!8] (O) -- (W) -- (V1W) -- (V1) -- cycle;
    \fill[black!20] (V1) -- (V1W) -- (SW) -- (S) -- cycle;
    \draw (V1) -- (V1W);
    \draw (S) -- (SW) node[midway,above=5pt] {\(w\)};
    \draw (W) -- (V1W) -- (SW);
    \draw[vector] (O) -- (W) node[midway,below=5pt] {\(w\)};
    \draw[vector] (O) -- (V1) node[midway,left=5pt] {\(v_1\)};
    \draw[vector] (V1) -- (S) node[midway,left=5pt] {\(v_2\)};
    \draw[vector,dashed] (O) -- (S);
    \draw[vector,dashed] (W) -- (SW)
      node[midway,right=5pt] {\(v_1+v_2\)};
  \end{scope}
  \begin{scope}[xshift=6.1cm]
    % Head-to-tail construction from PDF p. 8. Each right vertex is the
    % corresponding left vertex translated by the same base vector w.
    \coordinate (r-origin) at (0,0);
    \coordinate (r-base) at (2.4,0);
    \coordinate (r-top) at (0.7,2.1);
    \coordinate (r-v2) at (-1.3,-1.05);
    \coordinate (r-top-base) at ($(r-top)+(r-base)$);
    \coordinate (r-sum) at ($(r-top)+(r-v2)$);
    \coordinate (r-sum-base) at ($(r-sum)+(r-base)$);
    \fill[black!8] (r-origin) -- (r-base) -- (r-top-base) -- (r-top) -- cycle;
    \fill[black!20] (r-sum) -- (r-sum-base) -- (r-top-base) -- (r-top) -- cycle;
    \draw[vector] (r-origin) -- (r-base) node[midway,below=5pt] {\(w\)};
    \draw[vector] (r-origin) -- (r-top) node[pos=0.75,right=5pt] {\(v_1\)};
    \draw[vector] (r-top) -- (r-sum) node[midway,left=5pt] {\(v_2\)};
    \draw (r-top) -- (r-top-base) node[midway,above=5pt] {\(w\)};
    \draw (r-base) -- (r-top-base) -- (r-sum-base) -- (r-sum);
    \draw[vector,dashed] (r-origin) -- (r-sum)
      node[midway,left=5pt] {\(v_1+v_2\)};
    \draw[vector,dashed] (r-base) -- (r-sum-base);
  \end{scope}
\end{tikzpicture}

  \caption[Addition of oriented areas with a common base.]{Addition with a common base \(w\), shown schematically.
    The two input vectors lie on the same side of the base on the left,
    and on opposite sides on the right. Dashed sides represent \(v_1+v_2\).
    The signed areas add in both cases.}
  \label{fig:linear-algebra-i:determinant-area-addition}
\end{figure}

Rescaling a vector multiplies the oriented area by the same scalar;
reflection changes its sign. These operations appear in
\autoref{fig:linear-algebra-i:determinant-area-scaling}.
The two inputs are linear separately. Doubling both inputs therefore
multiplies the result by \(2\) twice,
\[
  A(2v,2w)=4A(v,w).
\]

\begin{figure}[htbp]
  \centering
  % Supplied October 1 PDF, pp. 9--10; schematic scaling and reflection.
\begin{tikzpicture}[
  x=1cm,y=1cm,font=\footnotesize,line width=0.5pt,
  line cap=round,line join=round,
  vector/.style={-{Latex[length=1.6mm,width=1.1mm]}}
]
  \begin{scope}
    \coordinate (O) at (0,0);
    \coordinate (W) at (2,0);
    \coordinate (V) at (0.5,1.1);
    \coordinate (S) at ($(V)+(W)$);
    \fill[black!8] (O) -- (W) -- (S) -- (V) -- cycle;
    \draw (V) -- (S) -- (W);
    \draw[vector] (O) -- (W) node[midway,below=5pt] {\(w\)};
    \draw[vector] (O) -- (V) node[midway,left=5pt] {\(v\)};
  \end{scope}
  \begin{scope}[xshift=3.8cm]
    \coordinate (O) at (0,0);
    \coordinate (W) at (2,0);
    \coordinate (V) at (0.5,1.1);
    \coordinate (CV) at ($2*(V)$);
    \coordinate (S) at ($(CV)+(W)$);
    \fill[black!8] (O) -- (W) -- (S) -- (CV) -- cycle;
    \draw (CV) -- (S) -- (W);
    \draw[vector] (O) -- (W) node[midway,below=5pt] {\(w\)};
    \draw[vector] (O) -- (CV) node[midway,left=5pt] {\(cv\)};
  \end{scope}
  \begin{scope}[xshift=8.2cm]
    \coordinate (O) at (0,0);
    \coordinate (W) at (2,0);
    \coordinate (V) at (0.5,1.1);
    \coordinate (NV) at ($-1*(V)$);
    \coordinate (S) at ($(NV)+(W)$);
    \fill[black!8] (O) -- (W) -- (S) -- (NV) -- cycle;
    \draw (NV) -- (S) -- (W);
    \draw[vector] (O) -- (W) node[midway,above=5pt] {\(w\)};
    \draw[vector] (O) -- (NV) node[midway,left=5pt] {\(-v\)};
  \end{scope}
\end{tikzpicture}

  \caption[Rescaling and reflection of oriented area.]{Rescaling and reflection of a parallelogram, shown schematically
    with \(c>1\). The identities are \(A(cv,w)=cA(v,w)\) and
    \(A(-v,w)=-A(v,w)\).}
  \label{fig:linear-algebra-i:determinant-area-scaling}
\end{figure}

\begin{proposition}[Further properties of oriented area]
  \label{thm:linear-algebra-i:oriented-area-properties}
  The oriented area function vanishes on equal inputs and is antisymmetric.
  For all \(v,w\in V\),
  \[
    A(v,v)=0,\qquad A(w,v)=-A(v,w).
  \]
\end{proposition}

Equal-input vanishing follows from the degenerate parallelogram described
above. For a bilinear function on a real vector space, this property implies
antisymmetry. These are additional properties of oriented area, rather than
conditions in the definition of bilinearity.
% PDF p. 12 and IMG_9207 state the implication without a proof.
% Lucas's October 6 proof is shared in the following exercise's solution.
The area function is an antisymmetric bilinear tensor.

The following exercise asks for the antisymmetry argument.
% Source: MA-GY7033 Fall 2026 Homework 5.pdf, p. 1, Problem 1.
\begin{exercise}
  \label{ex:hw05:bilinear-antisymmetry}
  \solutionlink{hw05:bilinear-antisymmetry}
  Let \(A:V\times V\to\R\) be a bilinear function on a real vector
  space \(V\) such that \(A(v,v)=0\) for every \(v\in V\).
  Show \(A\) is antisymmetric.
\end{exercise}


% Keep the coordinate derivation, its explanation, and its margin diagram together.
\newpage
\subsection{The determinant of a two-by-two matrix}

Write \(v=ae_1+ce_2\) and \(w=be_1+de_2\).
\autoref{fig:linear-algebra-i:determinant-two-by-two-area} shows these
coordinates and the parallelogram spanned by the two vectors.

\begin{marginfigure}
  \centering
  % Author-requested visualization of the existing two-by-two calculation,
% October 6, 2026. Coordinates illustrate one positive orientation only.
\begin{tikzpicture}[
  x=0.85cm,y=0.85cm,font=\footnotesize,line width=0.5pt,
  line cap=round,line join=round,
  vector/.style={-{Latex[length=1.6mm,width=1.1mm]}},
  axis/.style={vector,black!55,line width=0.4pt},
  guide/.style={black!50,line width=0.35pt,dashed}
]
  \coordinate (O) at (0,0);
  \coordinate (V) at (2.25,0.6);
  \coordinate (W) at (0.8,2.1);
  \coordinate (S) at ($(V)+(W)$);
  \coordinate (VX) at (V |- O);
  \coordinate (VY) at (V -| O);
  \coordinate (WX) at (W |- O);
  \coordinate (WY) at (W -| O);

  \fill[black!8] (O) -- (V) -- (S) -- (W) -- cycle;
  \draw[axis] (-0.15,0) -- (3.55,0) node[right=4pt] {\(e_1\)};
  \draw[axis] (0,-0.15) -- (0,3.2) node[above=4pt] {\(e_2\)};
  \draw[guide] (VX) -- (V) -- (VY);
  \draw[guide] (WX) -- (W) -- (WY);
  \draw (V) -- (S) -- (W);
  \draw[vector] (O) -- (V);
  \draw[vector] (O) -- (W);

  \draw (VX) ++(0,-0.045) -- ++(0,0.09);
  \draw (WX) ++(0,-0.045) -- ++(0,0.09);
  \draw (VY) ++(-0.045,0) -- ++(0.09,0);
  \draw (WY) ++(-0.045,0) -- ++(0.09,0);
  \node[below=5pt] at (VX) {\(a\)};
  \node[below=5pt] at (WX) {\(b\)};
  \node[left=5pt] at (VY) {\(c\)};
  \node[left=5pt] at (WY) {\(d\)};
  \node[right=4pt] at (V) {\(v\)};
  \node[above left=3pt] at (W) {\(w\)};
  \node[above=5pt] at (S) {\(v+w\)};
  \node[below left=3pt] at (O) {\(0\)};
\end{tikzpicture}

  \caption[Two-by-two determinant as oriented area.]{Basis-coordinate
    schematic for \(v=ae_1+ce_2\) and \(w=be_1+de_2\). The shaded
    parallelogram has oriented area \(ad-bc\), positive in the pictured
    configuration.}
  \label{fig:linear-algebra-i:determinant-two-by-two-area}
\end{marginfigure}

The equal-input property gives \(A(e_1,e_1)=A(e_2,e_2)=0\).
Bilinearity gives
\[
  \begin{aligned}
    A(v,w)
      &=A(ae_1,be_1)+A(ce_2,be_1)\\
      &\quad+A(ae_1,de_2)+A(ce_2,de_2)\\
      &=abA(e_1,e_1)+bcA(e_2,e_1)\\
      &\quad+adA(e_1,e_2)+cdA(e_2,e_2)\\
      &=bcA(e_2,e_1)+adA(e_1,e_2)\\
      &=ad-bc.
  \end{aligned}
\]
The terms \(abA(e_1,e_1)\) and \(cdA(e_2,e_2)\) vanish because their
inputs repeat. Antisymmetry gives
\(A(e_2,e_1)=-A(e_1,e_2)\), and the normalization is
\(A(e_1,e_2)=1\).

\begin{definition}
  The determinant of a two-by-two matrix is
  \[
    \det\begin{bmatrix}a&b\\c&d\end{bmatrix}=ad-bc.
  \]
\end{definition}

In column notation,
\[
  [v\ w]=[e_1\ e_2]\begin{bmatrix}a&b\\c&d\end{bmatrix}.
\]
Thus the determinant is the oriented area of the parallelogram whose
generating vectors are the columns, measured in the normalized basis.
Before imposing the normalization, the coordinate identity reads
\[
  A(v,w)=A(e_1,e_2)(ad-bc).
\]

\subsection{Oriented volume}

In a three-dimensional real vector space, three independent vectors
\(a,b,c\) span the parallelepiped
\[
  P(a,b,c)=\{sa+tb+uc\mid 0\leq s,t,u\leq1\}.
\]
Fix an ordered basis \((e_1,e_2,e_3)\) whose unit parallelepiped has
volume \(1\). Suppose \(a,b\) lie in the plane spanned by \(e_1,e_2\),
and write
\[
  c=ue_1+ve_2+he_3.
\]
Only the \(e_3\)-component takes \(c\) off the base plane. Its coefficient
\(h\) is the height in these coordinates, as shown in
\autoref{fig:linear-algebra-i:determinant-parallelepiped}.
The ordinary volume is
\[
  \operatorname{vol}P(a,b,c)=|h|\,|A(a,b)|.
\]

\begin{figure}[htbp]
  \centering
  % Supplied October 1 PDF, pp. 15--17, and IMG_9208; schematic projection.
\begin{tikzpicture}[
  x=1cm,y=1cm,font=\footnotesize,line width=0.5pt,
  line cap=round,line join=round,
  vector/.style={-{Latex[length=1.6mm,width=1.1mm]}},
  guide/.style={black!55,dashed}
]
  \coordinate (O) at (0,0);
  \coordinate (A) at (-0.95,-0.65);
  \coordinate (B) at (3.1,0);
  \coordinate (C) at (0.65,2.3);
  \coordinate (AB) at ($(A)+(B)$);
  \coordinate (AC) at ($(A)+(C)$);
  \coordinate (BC) at ($(B)+(C)$);
  \coordinate (ABC) at ($(A)+(B)+(C)$);
  \coordinate (E3) at (0,0.85);
  \coordinate (H) at (O |- C);
  \fill[black!8] (O) -- (B) -- (AB) -- (A) -- cycle;
  \fill[black!4] (C) -- (BC) -- (ABC) -- (AC) -- cycle;
  \draw (A) -- (AB) -- (B);
  \draw (A) -- (AC) -- (ABC) -- (AB);
  \draw (B) -- (BC) -- (ABC);
  \draw (AC) -- (C) -- (BC);
  \draw[guide] (O) -- (H) -- (C);
  \draw[vector] (O) -- (A) node[midway,left=16pt] {\(a\)};
  \draw[vector] (O) -- (B) node[midway,below=5pt] {\(b\)};
  \draw[vector] (O) -- (C) node[midway,right=5pt] {\(c\)};
  \draw[vector,black!55] (O) -- (E3) node[pos=0.6,left=5pt] {\(e_3\)};
  \node[above left=5pt] at (H) {\(he_3\)};
  \node[below=5pt] at (O) {\(0\)};
\end{tikzpicture}

  \caption[Parallelepiped and its height coefficient.]{A parallelepiped with base spanned by \(a,b\), shown
    schematically. The height coefficient is the \(e_3\)-component of
    \(c\); the drawing does not specify a Euclidean metric.}
  \label{fig:linear-algebra-i:determinant-parallelepiped}
\end{figure}

The oriented volume \(\mathcal V(a,b,c)\) is normalized by
\[
  \mathcal V(e_1,e_2,e_3)=1,
  \qquad |\mathcal V(a,b,c)|=\operatorname{vol}P(a,b,c).
\]
It is antisymmetric and linear in each of its three inputs. In the
base-plane configuration,
\[
  \mathcal V(a,b,c)=hA(a,b).
\]
This formula has no absolute values.

If \([v_1\ v_2\ v_3]=[e_1\ e_2\ e_3]M=EM\), then the determinant
of the three-by-three matrix \(M\) is characterized by
\[
  \mathcal V(v_1,v_2,v_3)=\mathcal V(e_1,e_2,e_3)\det M.
\]
With the chosen normalization, \(\det M=\mathcal V(v_1,v_2,v_3)\).
The same algebraic construction extends to \(n\) inputs through
multilinear functions.
